\documentclass[12pt]{article}
\usepackage{hyperref}

\usepackage{graphicx}

\usepackage{mathtools}
\DeclarePairedDelimiter\ceil{\lceil}{\rceil}
\DeclarePairedDelimiter\floor{\lfloor}{\rfloor}


\usepackage{amsmath}
\usepackage{amsfonts}
\usepackage{amssymb}

\usepackage{bm}

\usepackage{diagbox}

\usepackage{xcolor}

\DeclareMathOperator{\lcm}{lcm}


\usepackage{listings}
\lstset{language=C++,
                basicstyle=\ttfamily,
                keywordstyle=\color{blue}\ttfamily,
                stringstyle=\color{red}\ttfamily,
                commentstyle=\color{green}\ttfamily,
                morecomment=[l][\color{magenta}]{\#}
}

\usepackage{amsthm}
\newtheorem{theorem}{Theorem}


\usepackage{breqn}
\usepackage[]{algorithm2e}
\usepackage{physics}

\newcommand\fnurl[2]{%
  \href{#2}{#1}\footnote{\url{#2}}%
}

\usepackage{cancel}

\newcommand{\Four}{\mathcal{F}}
\renewcommand{\(}{\left(}
\renewcommand{\)}{\right)}
%% \renewcommand{\{}{\left{}
%% \renewcommand{\}}{\right}}


\title{Notes on Strides}
\author{The rocFFT Team}



\def\no#1{\nobreak{#1}} % stop equation breaking in dmath

\begin{document}
\maketitle

\section{Introduction}

rocFFT buffers are organized as $n$-dimensional arrays with
non-negative strides.  We currently support 3 dimensions with one
batch dimension, but we are interested in increasing this, with at
least 2 batch dimensions.

For a length $\bm\ell=\(\ell_0,\dots,\ell_{N-1}\)$, index
$\bm{a}=\(a_0, \dots, a_{N-1}\)$ has elements $a_i$ with
$0\leq{}a_i<\ell_i$. Given the stride $\bm{s}=\(s_0,\dots,s_{N-1}\)$,
index $\bm{a}$ refers to memory at pointer-offset
\begin{dmath}
  \bm{a}\cdot{}\bm{s}=\sum_{i=0}^{N-1}s_ia_i.
\end{dmath}
We require that $\bm{a}\cdot{}\bm{s}$ be injective; that is, no two
indices refer to the same pointer offset.

\section{Stride/length validity}

Not all combinations of strides and lengths produce are valid.  For
example, 1D transforms on non-unit length must have stride $>0$;
otherwise, index 0 and index 1 would both refer to data at offset 0,
and we would get a race condition when reading/writing to that
address.

The situation in higher dimensions is somewhat more complicated.  If
there existss two admissible indices $\bm{a}$ and $\bm{b}$ such that
\begin{dmath}
  \bm{a}\cdot\bm{p} = \bm{b}\cdot\bm{p}
\end{dmath}
then the stride/length combination is invalid.  When do such
combinations exist, and how can we detect them?

\begin{figure}[htbp]
  \begin{center}
    \begin{tabular}{|l|lllll|}
      \hline
      \diagbox{5}{3} & 0 & 1 & 2 & 3 & 4\\
      \hline
      0 & 0             & 3  & 6  & 9  & 12 \\
      1 & 5             & 8  & 11 & 14 & 17            \\
      2 & 10            & 13 & 16 & 19 & 22            \\
      3 & 15 & 18 & 21 & 24 & 27            \\
      \hline
    \end{tabular}
  \end{center}
  \caption{Pointer offsets for strides $(3,5)$ and lengths $(5,4)$.}
  \label{tab35_len54}
\end{figure}


Take length $(\ell_0,\ell_1)$ with stride $\bm{s} = (s_0,s_1)$.  The
lowest address that has an index collission is at $\lcm(s_0,s_1)$.
For example, with strides $\bm{s}=(3,5)$, the first few pointer
offsets are given in Figure~\ref{tab35_len64}.
\begin{figure}[htbp]
  \begin{center}
    \begin{tabular}{|l|llllll|}
      \hline
      \diagbox{5}{3} & 0 & 1 & 2 & 3 & 4 & 5\\
      \hline
      0 & 0             & 3  & 6  & 9  & 12 & \color{red}15 \\
      1 & 5             & 8  & 11 & 14 & 17 & 20            \\
      2 & 10            & 13 & 16 & 19 & 22 & 25            \\
      3 & \color{red}15 & 18 & 21 & 24 & 27 & 30            \\
      \hline
    \end{tabular}
  \end{center}
  \caption{Pointer offsets for strides $(3,5)$ and lengths $(6,4,)$.}
  \label{tab35_len64}
\end{figure}
In fact, it's easy to see that collisions will only occur at multiples
of $\lcm(s_0,s_1)$.
\begin{theorem}
  Given the of positive strides $s_0$ and $s_1$, the first invalid
  data point occurs at address $\lcm{(s_0,s_1)}$.
\end{theorem}
\begin{proof}
  Suppose that $\gcd{s_0,s_1}=1$.  We know that $(s_1,0)$ and
  $(0,s_0)$ give the address $s_0 s_1=\lcm(s_0,s_1)$.  Thus, the first
  invalid address is at most $\lcm(s_0,s_1)$.

  Suppose that there exists valid index $(a_0,a_1)$ and
  $c<\lcm(s_0,s_1)$ such that
  \begin{dmath}
    s_0 a_0 + s_1 a_1 = c.
  \end{dmath}
  By Bézout's identity, we can find $x,y\in\mathbb{Z}$ such that
  \begin{dmath}
    x s_0 + y s_1 = 1
  \end{dmath}
  since $\gcd(s_0,s_1)=1$.  Then, $(c x, c y)$ gives address $c$.  We
  are guaranteed multiple solutions over $\mathbb{Z}^2$, namely
\begin{dmath}
  (cx + s_1t, cy - s_0t)
\end{dmath}
for arbitrary $t\in\mathbb{Z}$.  However, since $c<s_0s_1$, if
$t\neq{}0$, then either $cx + s_1t$ is negative, or $cy - s_0t$ is, and
there is only one solution in $\mathbb{N}^2$. Thus, $c=\lcm(s_0,s_1)$.

If $\gcd{s_0,s_1}=g>1$, then, again by Bézout, all memory addresses
are multiple of $g$.  Noting that
$$
s_0 a_0 + s_1 a_1 = c \implies
\frac{s_0}{g} a_0 + \frac{s_1}{g} a_1 = \frac{c}{g},
$$
and that
\begin{dmath}
  \lcm\(\frac{s_0}{g}, \frac{s_1}{g}\) = 1
\end{dmath},
the $\gcd{s_0,s_1}>1$ case reduces to the $\gcd{s_0,s_1}>1$ case.
\end{proof}
Unlike the Bézout's identity, not all redundant linear combinations
are multiples of the $\lcm$.  This is easy to see with a counter-example.

Consider strides $(2,3)$, which has $\lcm=6$.  Index $(1,2)$ gives
memory position $1 \times{} 2 + 2 \times{} 3 = 8$, which is also attained by
$(4,0)\rightarrow{}4\times{}2 + 0\times{}3 = 8$.



The simplest way for an index to attain the $\lcm$ is for one index to
be zero: we are guaranteed that
\begin{dmath}
  \(\frac{\lcm(s_0,s_1))}{s_0}, 0\) \no\cdot (s_0, s_1) 
  =  \(0, \frac{\lcm(s_0,s_1))}{s_1}\) \no\cdot (s_0, s_1) 
  = \lcm(s_0,s_1).
\end{dmath}
Thus, we require that at least one of
\begin{eqnarray}
\ell_0<\frac{\lcm(s_0,s_1))}{s_0}\\
\ell_1<\frac{\lcm(s_0,s_1))}{s_1}
\end{eqnarray}
be false.


In the remaining cases, where at one length is long enough to reach
the lcm, but not all are, and the lcm is attainable, then we can
brute-force the remaining values.  In particular, we need to check,
for $a_0 = 1, \dots, \lcm(s_0,s_1)/s_0 - 1$, if we can find $a_1$ such
that
\begin{dmath}
  a_1 = \frac{\lcm(s_0,s_1) - a_0 s_0}{s_1}.
\end{dmath}
In particular, it suffices to check whether 
\begin{dmath}
  s_1  \vert \(\lcm(s_0,s_1) - a_0 s_0\)
\end{dmath}
for the above values of $a_0$. 



\section{Buffer size computation}

ways to calculate:
sum
last index
max of extrema


\section{negative indices}

\end{document}
